\chapter{Canonical Data Schema}
\label{chap:schema}\label{sim:chap:schema}

\section{Layout of \texttt{src/training/pipeline/}}

\begin{table}[H]
\centering
\caption{Pipeline Module Layout}
\label{tab:module-layout}
\begin{tabularx}{\textwidth}{lX}
\toprule
\textbf{Path} & \textbf{Purpose} \\
\midrule
\texttt{simula\_config.py} & Configuration dataclasses for all components \\
\texttt{simula\_llm\_client.py} & OpenAI-compatible async LLM client \\
\texttt{hana\_schema\_extractor.py} & HANA Cloud schema extraction via MCP \\
\texttt{simula\_taxonomy\_builder.py} & Algorithm 1: Taxonomy generation \\
\texttt{simula\_data\_generator.py} & Algorithm 2: Training data synthesis \\
\texttt{schema\_registry.py} & In-memory schema registry \\
\bottomrule
\end{tabularx}
\end{table}

\section{Training Example Record}

Each generated training example is a JSON object with the following structure:

\begin{lstlisting}[language=jsonx,caption={TrainingExample dataclass (simula\_data\_generator.py)}]
@dataclass
class TrainingExample:
    """A single training example for Text-to-SQL."""
    id: str                    # Unique identifier
    question: str              # Natural language question
    sql: str                   # Target SQL query
    schema_context: str        # Table/column descriptions
    taxonomy_path: list[str]   # Factor nodes used
    complexity_score: float    # 0.0 to 1.0
    meta_prompt_id: str        # Reference to meta-prompt record
    audience: str              # human, agent, or dual
    is_complexified: bool      # Whether complexification applied
    critic_passed: bool        # Whether double-critic passed
    generation_metadata: dict  # Timestamps, model, etc.
\end{lstlisting}

\section{Schema Registry}

The \texttt{SchemaRegistry} holds table and column metadata extracted from HANA:

\begin{lstlisting}[language=jsonx,caption={Schema registry structures (schema\_registry.py)}]
@dataclass
class Column:
    name: str
    data_type: str
    length: int | None
    nullable: bool
    description: str = ""

@dataclass
class TableSchema:
    schema_name: str
    table_name: str
    columns: list[Column]
    primary_key: list[str] = field(default_factory=list)
    description: str = ""
\end{lstlisting}

\section{Taxonomy Structure}

Taxonomies are hierarchical trees where each node represents a factor of variation:

\begin{lstlisting}[language=jsonx,caption={Taxonomy structures (simula\_taxonomy\_builder.py)}]
@dataclass
class TaxonomyNode:
    id: str
    name: str
    description: str
    level: int
    parent_id: str | None
    children: list["TaxonomyNode"]
    examples: list[str]  # Sample values

@dataclass
class Factor:
    id: str
    name: str
    description: str
    taxonomy: "Taxonomy"

@dataclass
class Taxonomy:
    id: str
    factor: str
    root: TaxonomyNode
    depth: int
    total_nodes: int
\end{lstlisting}

\section{Enumerations}

\begin{table}[H]
\centering
\caption{Key Enumerations}
\label{tab:enums}
\begin{tabular}{lll}
\toprule
\textbf{Enum} & \textbf{Values} & \textbf{Description} \\
\midrule
\texttt{LLMProvider} & \textsc{vllm}, \textsc{openai}, \textsc{anthropic} & LLM backend type \\
\texttt{SamplingStrategy} & \textsc{uniform}, \textsc{weighted}, \textsc{coverage} & How to sample taxonomy nodes \\
\bottomrule
\end{tabular}
\end{table}

\section{Unified Schema Registry}

As of 18 April 2026, all schemas in this specification are published in the
unified schema registry at \path{docs/schema/}. Migration status is
\emph{complete} (see \path{docs/schema/registry.json} field
\texttt{migration.status}). This consolidation addresses cross-document schema
divergence identified in the specification review.

\begin{description}
  \item[Registry index.] \path{docs/schema/registry.json} --- Master index of
        all schemas across domains.
  \item[Common types.] \path{docs/schema/common/base-types.schema.json} ---
        Shared definitions (\texttt{Metadata}, \texttt{Currency}, etc.)
        referenced via JSON Schema \texttt{\$ref}.
  \item[Consolidated enums.] \path{docs/schema/common/enums.yaml} --- Single
        source of truth for enumeration values used by all four specifications.
  \item[Simula schemas.] \path{docs/schema/simula/} --- Domain-specific
        schemas: \texttt{taxonomy.schema.json}, \texttt{training-example.schema.json},
        \texttt{hana-schema-entry.schema.json}, \texttt{config.schema.json}, and
        \texttt{meta-prompt.schema.json}. All marked \texttt{status: stable}.
\end{description}

\noindent
The registry uses JSON Schema Draft-07 with the canonical \texttt{\$id} prefix
\url{https://sap-oss.github.io/schema/}. Validation:

\begin{lstlisting}[language=bash]
python3 docs/schema/validate.py --registry-check
python3 docs/schema/validate.py --domain simula \
    --schema training-example.schema.json --file example.json
\end{lstlisting}

\noindent
See \path{docs/schema/README.md} for Python usage examples; the CI pipeline
executes \texttt{--registry-check} on every pull request that touches
\path{docs/schema/} or any spec chapter that references these schemas.

% =============================================================================
% CHAPTER 3: HANA SCHEMA EXTRACTION PIPELINE
% =============================================================================
\newpage
